% ============================================================
%  MEMO EJECUTIVO - CONCEPT BPO
%  Compilar con: lualatex o xelatex memo_ejecutivo_concept_bpo.tex
% ============================================================
\documentclass[10pt,a4paper]{article}

\usepackage{fontspec}
\usepackage[spanish]{babel}
\usepackage[top=1.5cm,bottom=1.5cm,left=2.2cm,right=2.2cm]{geometry}
\usepackage{xcolor}
\usepackage{tabularx}
\usepackage{parskip}
\usepackage{microtype}
\usepackage{enumitem}
\usepackage{mdframed}
\usepackage{titlesec}
\usepackage{fancyhdr}

\definecolor{cbpoBlue}{HTML}{1E3A5F}
\definecolor{cbpoAccent}{HTML}{2563EB}
\definecolor{cbpoRed}{HTML}{DC2626}
\definecolor{cbpoGray}{HTML}{6B7280}
\definecolor{cbpoLight}{HTML}{EFF6FF}

\pagestyle{fancy}
\fancyhf{}
\renewcommand{\headrulewidth}{0pt}
\fancyhead[L]{\small\color{cbpoGray}\textsf{CONFIDENCIAL -- Uso interno}}
\fancyhead[R]{\small\color{cbpoGray}\textsf{Concept BPO}}
\fancyfoot[C]{\small\color{cbpoGray}\thepage}

\titleformat{\section}{\sffamily\bfseries\color{cbpoBlue}\large}{}{0em}{}[\color{cbpoAccent}\rule{\linewidth}{0.8pt}]
\titlespacing{\section}{0pt}{6pt}{2pt}

\newmdenv[
  linecolor=cbpoAccent,linewidth=2pt,backgroundcolor=cbpoLight,
  leftmargin=0pt,rightmargin=0pt,
  innerleftmargin=10pt,innerrightmargin=10pt,
  innertopmargin=8pt,innerbottommargin=8pt,
  skipabove=8pt,skipbelow=4pt
]{hallazgobox}

\begin{document}
\thispagestyle{fancy}

{\color{cbpoBlue}
\noindent\rule{\linewidth}{2pt}\\[4pt]
\begin{tabularx}{\linewidth}{@{}lX@{}}
  \textbf{\textsf{\Large MEMO EJECUTIVO}} & \\[4pt]
  \textsf{\textbf{Para:}}          & Director de Operaciones, Concept BPO \\
  \textsf{\textbf{De:}}            & Equipo de Ciencia de Datos \\
  \textsf{\textbf{Fecha:}}         & \today \\
  \textsf{\textbf{Asunto:}}        & Estrategia de reducción de abandono y optimización operativa \\
  \textsf{\textbf{Clasificación:}} & \textcolor{cbpoRed}{\textbf{Confidencial}} \\
\end{tabularx}\\[4pt]
\noindent\rule{\linewidth}{2pt}
}

\section{Hallazgo Principal}

\begin{hallazgobox}
\textbf{El 15.4\,\% de las llamadas se pierden por abandono, concentrándose en picos de congestión donde la disponibilidad de agentes cae por debajo del nivel crítico.}
Esto representa una pérdida de \textbf{2\,300 interacciones mensuales}, lo que no solo afecta el CSAT, sino que genera un efecto de "bola de nieve": el cliente que cuelga vuelve a llamar minutos después, saturando aún más la cola de espera.
\end{hallazgobox}

El análisis de 15\,000 registros identifica que el abandono no es aleatorio, sino altamente predecible bajo tres condiciones:
\begin{itemize}[leftmargin=1.4em,itemsep=1pt]
  \item \textbf{Umbral Crítico de Disponibilidad:} El modelo predictivo confirma que la variable más influyente es el número de agentes disponibles en el instante de entrada. Existe un punto de quiebre donde la probabilidad de abandono se duplica instantáneamente.
  \item \textbf{Sensibilidad por Sector:} Los clientes de \textbf{Retail} son los más sensibles; tienen los tiempos de espera más altos en la mañana (102\,s) y abandonan con mayor rapidez que los de Banca.
  \item \textbf{Fugas por Re-llamada:} La baja resolución en primer contacto (FCR) de 19 agentes específicos está re-inyectando volumen innecesario a la cola, agravando la espera en horas pico.
\end{itemize}

\section{Recomendación y Estimación de Impacto}

Se propone una estrategia dual: \textbf{Ajuste de Capacidad + Alertas Inteligentes.}

\begin{enumerate}[leftmargin=1.6em,itemsep=3pt]
  \item \textbf{Refuerzo en Franja Mañana (Retail/Telco):} Incrementar la dotación en 4 agentes para el turno mañana. Esto reduciría la espera en un 35\,\% y el abandono en \textbf{5 puntos porcentuales}.
  \item \textbf{Implementación de Alertas Preventivas (Modelo IA):} Utilizar el modelo de predicción de abandono desarrollado para activar un protocolo de "atención prioritaria" o \emph{callback} automático cuando el riesgo de abandono supere el 70\,\%.
\end{enumerate}

\medskip
\noindent\textbf{Impacto esperado:} Recuperación de $\sim$850 llamadas mensuales, mejora del CSAT en 0.3 puntos y reducción de la presión operativa por re-llamadas evitadas.

\section{Limitaciones del Análisis}

\begin{enumerate}[leftmargin=1.6em,itemsep=1pt]
  \item \textbf{Sesgo en Satisfacción:} Solo disponemos de CSAT para el 51\,\% de las llamadas (las completadas), por lo que la visión de la experiencia del cliente está incompleta.
  \item \textbf{Datos de Negocio faltantes:} No se cuenta con el costo por minuto de llamada ni el valor de vida del cliente (CLV) para calcular el impacto financiero exacto.
  \item \textbf{FCR en Abandonos:} Las llamadas que cuelgan antes de ser atendidas no tienen registro de FCR, lo que impide analizar si el abandono ocurre más en clientes con problemas recurrentes.
\end{enumerate}

\vspace{12pt}
\noindent\small\color{cbpoGray}
\textit{Análisis basado en modelos de aprendizaje supervisado (Random Forest) con balanceo de clases. El soporte técnico y notebook detallado se adjuntan al presente.}

\end{document}
